\documentclass[../main.tex]{subfiles}

\begin{document}

\chapter{Rubrik Penilaian Lengkap Tertaut CPMK}

Lampiran ini menjabarkan bobot pada RPS menjadi rubrik yang dapat langsung
dipakai untuk menilai setiap instrumen. Seluruh penilaian bermuara pada dua
CPMK mata kuliah ini.

\section{CPMK dan Sub-CPMK}

\begin{table}[H]
\centering
\caption{Capaian pembelajaran mata kuliah dan turunannya}
\label{tab:cpmk-subcpmk}
\footnotesize
\begin{tabularx}{\linewidth}{|>{\centering\arraybackslash}p{1.6cm}|>{\raggedright\arraybackslash}X|>{\raggedright\arraybackslash}p{3.4cm}|}
\hline
\textbf{CPMK} & \textbf{Rumusan} & \textbf{Sub-CPMK} \\
\hline
CPMK-1 & Mengenal beberapa jenis algoritma kriptografi klasik dan modern & 1.1, 1.2, 1.3, 1.4 \\
\hline
CPMK-2 & Mahasiswa mampu memahami konsep kriptografi secara umum dan urgensinya dalam dunia teknologi informasi & 2.1, 2.2, 2.3 \\
\hline
\end{tabularx}
\end{table}

\section{Rubrik Tugas (15\%)}

\begin{table}[H]
\centering
\caption{Rubrik penilaian tugas}
\label{tab:rubrik-tugas}
\footnotesize
\begin{tabularx}{\linewidth}{|>{\raggedright\arraybackslash}X|>{\centering\arraybackslash}p{1.6cm}|>{\raggedright\arraybackslash}X|}
\hline
\textbf{Kriteria} & \textbf{Bobot} & \textbf{Indikator} \\
\hline
Ketepatan jawaban latihan & 40\% & Hasil akhir benar untuk setiap butir \textit{Latihan dan Refleksi}. \\
\hline
Kelengkapan langkah & 30\% & Langkah perhitungan ditulis runtut dan dapat diikuti. \\
\hline
Ketepatan istilah dan rujukan & 15\% & Istilah teknis tepat dan rujukan disebut benar. \\
\hline
Ketepatan waktu & 15\% & Dikumpulkan sesuai tenggat. \\
\hline
\end{tabularx}
\end{table}

\section{Rubrik Praktikum (20\%)}

\begin{table}[H]
\centering
\caption{Rubrik penilaian praktikum}
\label{tab:rubrik-praktikum}
\footnotesize
\begin{tabularx}{\linewidth}{|>{\raggedright\arraybackslash}X|>{\centering\arraybackslash}p{1.6cm}|>{\raggedright\arraybackslash}X|}
\hline
\textbf{Kriteria} & \textbf{Bobot} & \textbf{Indikator} \\
\hline
Kesesuaian prosedur & 30\% & Langkah kerja pada \textit{Aktivitas Rekomendasi} dijalankan benar. \\
\hline
Ketepatan hasil & 30\% & Keluaran program atau perhitungan cocok dengan acuan. \\
\hline
Dokumentasi laporan & 25\% & Laporan memuat langkah, hasil, dan pembahasan. \\
\hline
Kerja sama tim & 15\% & Pembagian peran jelas dan kontribusi merata. \\
\hline
\end{tabularx}
\end{table}

\section{Rubrik Proyek (15\%)}

\begin{table}[H]
\centering
\caption{Rubrik penilaian proyek}
\label{tab:rubrik-proyek}
\footnotesize
\begin{tabularx}{\linewidth}{|>{\raggedright\arraybackslash}X|>{\centering\arraybackslash}p{1.6cm}|>{\raggedright\arraybackslash}X|}
\hline
\textbf{Kriteria} & \textbf{Bobot} & \textbf{Indikator} \\
\hline
Identifikasi masalah & 20\% & Kasus nyata dirumuskan dengan jelas dan relevan. \\
\hline
Ketepatan pemilihan algoritma & 25\% & Algoritma yang dipilih sesuai kebutuhan keamanan dan efisiensi. \\
\hline
Implementasi dan verifikasi & 25\% & Algoritma diterapkan dan hasilnya diverifikasi. \\
\hline
Laporan dan presentasi & 20\% & Laporan sistematis dan presentasi jelas. \\
\hline
Orisinalitas & 10\% & Ada analisis atau pengembangan sendiri. \\
\hline
\end{tabularx}
\end{table}

\section{Kisi-kisi UTS dan UAS}

\begin{table}[H]
\centering
\caption{Distribusi soal ujian menurut Sub-CPMK}
\label{tab:kisi-ujian}
\footnotesize
\begin{tabularx}{\linewidth}{|>{\centering\arraybackslash}p{1.6cm}|>{\raggedright\arraybackslash}p{1.6cm}|>{\raggedright\arraybackslash}X|>{\centering\arraybackslash}p{1.8cm}|}
\hline
\textbf{Ujian} & \textbf{Sub-CPMK} & \textbf{Cakupan materi} & \textbf{Bobot} \\
\hline
UTS & 2.1, 2.2 & Pengantar, urgensi, layanan keamanan, serangan & 30\% \\
\hline
UTS & 2.3 & Landasan matematika & 15\% \\
\hline
UTS & 1.1 & Cipher klasik dan kriptanalisis & 25\% \\
\hline
UTS & 1.2 & Stream cipher, block cipher, DES & 30\% \\
\hline
UAS & 1.2 & AES dan mode operasi & 25\% \\
\hline
UAS & 1.3 & RSA, Diffie--Hellman, ElGamal, ECC & 40\% \\
\hline
UAS & 1.4 & Hash, MAC, tanda tangan digital, PKI, SSL/TLS & 35\% \\
\hline
\end{tabularx}
\end{table}

Bobot pada setiap ujian dibulatkan terhadap 100\%; butir soal dapat mencakup
lebih dari satu Sub-CPMK.

\end{document}
